% Generated draft: requires mathematical and source review.
\documentclass[11pt]{article}
\usepackage[margin=25mm]{geometry}
\usepackage{fontspec}
\setmainfont{Latin Modern Roman}
\usepackage{amsmath,amssymb,mathtools,unicode-math}
\setmathfont{Latin Modern Math}
\usepackage{enumitem,needspace,xurl,hyperref}
\setlength{\parindent}{0pt}
\setlength{\parskip}{5pt}
\setlength{\emergencystretch}{3em}
\newcommand{\N}{\mathbb{N}}
\newcommand{\Z}{\mathbb{Z}}
\newcommand{\Q}{\mathbb{Q}}
\newcommand{\R}{\mathbb{R}}
\newcommand{\C}{\mathbb{C}}
\newcommand{\F}{\mathbb{F}}
\newcommand{\A}{\mathbb{A}}
\newcommand{\PP}{\mathbb{P}}
\newcommand{\E}{\mathbb{E}}
\newcommand{\eps}{\varepsilon}
\newcommand{\dd}{\,\mathrm{d}}
\begin{document}
\section*{UnsolvedMath ID: 30003498}
\textbf{Source problem code:} catalogue entry
\subsection{Definitions and mathematical statement}
The supplied record poses two definability questions, but does not specify a product formula or all definability conventions. The following makes a standard literal interpretation explicit; these conventions require source review.

Let \(\mathcal R=(\R;0,1,+,\cdot,<)\) be the ordered real field. A function on a definable domain is definable in a structure if its graph is definable in that structure. In the conditional formulation below, real parameters are allowed; the record does not specify whether parameters are permitted.

For an integer \(p\geq 0\), put
\[
E_p(w)=(1-w)\exp\!\left(\sum_{j=1}^{p}\frac{w^j}{j}\right),
\]
where the sum is zero when \(p=0\). Taking the listed sequence to consist of positive zeros, taking indices \(n=1,2,\ldots\), and using the least admissible genus, define, for each fixed real \(s\) with \(0<s<1\),
\[
p_s=\min\{p\in\Z_{\geq0}:s(p+1)>1\},
\qquad
P_s(z)=\prod_{n=1}^{\infty}E_{p_s}\!\left(\frac{z}{n^s}\right),
\quad z\in\C.
\]
Under these explicit conventions, the first question is:
\[
\text{For every }s\in(0,1),\text{ is }
\{(x,y)\in\R^2:x>0,\ y=P_s(x)\}
\text{ definable in }\mathcal R\text{?}
\]
This asks about each fixed \(s\); it does not impose uniform definability as \(s\) varies. Neither endpoint \(s=0\) nor \(s=1\) is included.

For each fixed real \(\alpha>1\), the analogous genus-zero, positive-zero convention gives
\[
Q_\alpha(z)=\prod_{n=1}^{\infty}\left(1-\frac{z}{\alpha^n}\right),
\qquad
D_\alpha=\alpha^{\Z}=\{\alpha^k:k\in\Z\}.
\]
If the intended added function is its restriction to the positive real axis, write
\[
\mathcal M_\alpha=
\bigl(\mathcal R;\operatorname{graph}(Q_\alpha|_{(0,\infty)})\bigr).
\]
For \(m\geq1\), let \(\operatorname{Def}_m(\mathcal M_\alpha)\) be the collection of subsets of \(\R^m\) definable in this structure, with the parameter convention above. The second question requests a description, for every \(\alpha>1\), of the definable sets in the expansion by the associated canonical product, beyond the set \(D_\alpha\), which the supplied statement describes as definable. It does not give a specific proposed classification or a precise criterion for ``beyond.''

The product formulas and the positive-axis domain for the second expansion are conditional specifications, not recovered source facts. The record does not settle the indexing, normalization, genus convention, or whether the second expansion uses a function on all of \(\R\). The intended ambient structure in the first question also needs confirmation: the literal pure-real-field wording and prescribed infinite positive zero sequence require reconciliation with the supplied assessment that the question is open. No alternative ambient structure or corrected target is asserted here.
\subsection{Short English statement}
For each \(0<s<1\), the record asks whether the canonical product associated with \((n^s)\), restricted to \((0,\infty)\), is definable in the real field. For each \(\alpha>1\), it also asks which sets become definable after adjoining the canonical product associated with \((\alpha^n)\), beyond \(\alpha^{\Z}\). The displayed products give an explicitly conditional interpretation; the intended first question, product conventions, and scope of the second classification need human source review. The supplied open status is dated curation, not an independently verified conclusion.
\subsection{Sources}
\begin{itemize}
\item Local UnsolvedMath registry, supplied complete source record: permanent ID \(30003498\), source problem code \texttt{catalogue entry}, title \emph{Definability of Canonical Products with Sparse Zero Sets}, in the collection \emph{Oberwolfach Reports Open Problems}. The record gives proposed year 2017 and names no individual proposer.
\item Supplied report citation: \emph{O-Minimality and its Applications to Number Theory and Analysis (2018)}, Oberwolfach Reports, pp.~1349--1420, DOI \texttt{10.4171/owr/2017/22}. \url{https://doi.org/10.4171/owr/2017/22}. The supplied literature-source label is ``Oberwolfach Reports 2018.'' The report is described as documenting mathematical workshops and their open-problem discussions.
\item Supplied literature source: Miller and Speissegger, ``canonical products.'' \url{https://arxiv.org/abs/1812.05547}.
\item Supplied literature source: Miller and Speissegger, ``published canonical-products results.'' \url{https://doi.org/10.4153/S0008439519000572}.
\item Literature assessment checked at \texttt{2026-08-21}: the supplied curation labels the problem open and states that canonical-product expansions have been analyzed for important parameter ranges, while the stated question for every \(0<s<1\) and the requested expansion remain open. This is supplied dated curation derived from the OWR material, not independently verified literature assessment; the status and cited claims require independent verification.
\end{itemize}
\end{document}
